
\medskip
\noindent
\SGARefTarget{sga7.sga7i.eix.n1-2}{1.2}. Apply VIII~\SGARefLink{sga7.sga7i.eviii.n7-1}{7.1}(b) to the pair $(A,A')$ and to the biextension
$W_K$ of $(A_K,A'_K)$ by $\mathbb G_{mK}$ in (\SGARefLink{sga7.sga7i.eix.n1-0-1}{1.0.1}).  We thereby
obtain a \underline{canonical pairing}
\begin{equation}\tag{\SGARefTarget{sga7.sga7i.eix.n1-2-1}{1.2.1}}
 \Phi_o\times\Phi'_o\longrightarrow(\mathbb Q/\mathbb Z)_k,
\end{equation}
that is, a pairing of $\pi_o$-modules
\begin{equation}\tag{\SGARefTarget{sga7.sga7i.eix.n1-2-2}{1.2.2}}
 \Phi_o(\overline k)\times\Phi'_o(\overline k)
 \longrightarrow\mathbb Q/\mathbb Z.
\end{equation}
Here the pairing appears as the obstruction to extending $W_K$ to a
biextension $W$ of $(A,A')$ by $\mathbb G_{mS}$; the meaning of this
point of view is explained in VIII~\SGARefLink{sga7.sga7i.eviii.n7-3-2}{7.3.2}.  We record at once the
following conjecture concerning this pairing.

\medskip
\noindent
\underline{Conjecture \SGARefTarget{sga7.sga7i.eix.n1-3}{1.3}.} \underline{The pairing} (\SGARefLink{sga7.sga7i.eix.n1-2-1}{1.2.1})
\underline{is a perfect duality}: \underline{it identifies each of the
component groups} $\Phi_o$ \underline{and} $\Phi'_o$
\underline{with the dual of the other, with values in}
$(\mathbb Q/\mathbb Z)_k$.

\medskip
\noindent
\SGARefTarget{sga7.sga7i.eix.n1-3-1}{1.3.1}. The formation of (\SGARefLink{sga7.sga7i.eix.n1-2-1}{1.2.1}) commutes with formally etale base
change (VIII~\SGARefLink{sga7.sga7i.eviii.n7-3-5-3}{7.3.5.3}), and the formation of Neron models has the
analogous property.  Consequently, to prove Conjecture~\SGARefLink{sga7.sga7i.eix.n1-3}{1.3} it would
suffice to treat the case in which $S$ is complete and has separably
closed residue field.  Artin and Mazur proved the conjecture when
$A_K$ is the Jacobian of a proper smooth curve over $K$, using their
general, unpublished theory of the self-duality of the relative
Jacobian (in the sense of derived categories) of a proper curve scheme
over $S$.  They also proved Conjecture~\SGARefLink{sga7.sga7i.eix.n1-3}{1.3} when the residue field is
\underline{finite}, using Tate duality [32].  We shall likewise see in
\SGARefLink{sga7.sga7i.eix.n11-4}{11.4} and \SGARefLink{sga7.sga7i.eix.n11-5}{11.5} (subject to compatibility checks) that \SGARefLink{sga7.sga7i.eix.n1-3}{1.3} holds on the
$\ell$-primary components for $\ell\ne p$, and that it holds in the
case of ``semistable reduction''.  These results make \SGARefLink{sga7.sga7i.eix.n1-3}{1.3} extremely
plausible.  Moreover, \SGARefLink{sga7.sga7i.eix.n1-3}{1.3} may also be regarded as a by-product of a
general duality theory for group schemes over fraction fields of
complete discrete valuation rings, formally analogous to SGA~5~I.
Such a theory remains heuristic at present; it would also contain
local class field theory in the form given to it by Serre [27].

\medskip
\noindent
\SGARefTarget{sga7.sga7i.eix.n1-3-2}{1.3.2}. For a given abelian scheme $A_K$, the validity of \SGARefLink{sga7.sga7i.eix.n1-3}{1.3} means that
the pairing (\SGARefLink{sga7.sga7i.eix.n1-2-1}{1.2.1}) is \underline{separating} on both the left and the
right.  By VIII~\SGARefLink{sga7.sga7i.eviii.n7-3-2}{7.3.2}, separation on the left says precisely that the
only open subgroup scheme $V$ of $A$ for which $W_K$ can be extended to
a biextension of $(V,A')$ by $\mathbb G_{mS}$ is the identity component
$A^o$ of $A$.  There is a symmetric interpretation of separation on
the right.

\medskip
\noindent
\SGARefTarget{sga7.sga7i.eix.n1-4}{1.4}. More generally, suppose that
\begin{equation}\tag{\SGARefTarget{sga7.sga7i.eix.n1-4-1}{1.4.1}}
 U\subset A\qquad\text{and}\qquad U'\subset A'
\end{equation}
are open subgroup schemes such that $W_K$ extends to a biextension of
$(U,U')$ by $\mathbb G_{mS}$.  By VIII~\SGARefLink{sga7.sga7i.eviii.n7-3-2}{7.3.2}, this says explicitly that
the subgroups
\begin{equation}\tag{\SGARefTarget{sga7.sga7i.eix.n1-4-2}{1.4.2}}
 F_o=U_o/U_o^o\subset\Phi_o,
 \qquad
 F'_o=U'_o/{U'}_o^o\subset\Phi'_o
\end{equation}
defined by $U$ and $U'$ are orthogonal under (\SGARefLink{sga7.sga7i.eix.n1-2-1}{1.2.1}).\footnote{The
printed formula, confirmed in a direct high-resolution check, reuses
$\Phi'_o$ for the first subgroup and writes $\Psi'_o\subset\Psi_o$ for
the second, colliding with the established notation for the component
group of $A'$.  The neutral local names $F_o,F'_o$ above make the
intended two subgroups unambiguous without changing the assertion.}
We denote this extension (unique up to unique isomorphism) of $W_K$ by
$W$.  It is determined up to unique isomorphism by its class
\begin{equation}\tag{\SGARefTarget{sga7.sga7i.eix.n1-4-3}{1.4.3}}
 w\in\mathrm{Biext}^1(U,U';\mathbb G_{mS}),
\end{equation}
which is itself characterized by the condition that its image in
$\mathrm{Biext}^1(A_K,A'_K;\mathbb G_{mK})$ be the class $w_K$ of
(\SGARefLink{sga7.sga7i.eix.n1-0-1}{1.0.1}).

For example, one may take for $U$ and $U'$ the identity components
$A^o$ and $A'{}^o$ of $A$ and $A'$.  One thereby obtains a
\underline{canonical biextension $W^o$, or $W$, of $(A^o,A'{}^o)$ by
$\mathbb G_{mS}$}, giving rise to a class
\begin{equation}\tag{\SGARefTarget{sga7.sga7i.eix.n1-4-4}{1.4.4}}
 w=w^o\in\mathrm{Biext}^1(A^o,A'{}^o;\mathbb G_{mS}),
\end{equation}
with which we shall work chiefly.  Evidently, the restrictions to
$(A^o,A'{}^o)$ of the classes $w$ in (\SGARefLink{sga7.sga7i.eix.n1-4-3}{1.4.3}), for the various pairs
$(U,U')$, are all equal to $w^o$ in (\SGARefLink{sga7.sga7i.eix.n1-4-4}{1.4.4}).

\medskip
\noindent
\underline{Proposition \SGARefTarget{sga7.sga7i.eix.n1-5}{1.5}.} \underline{Consider the special fiber
$W_s$ of the canonical biextension} $W$ \underline{of}
$(A^o,A'{}^o)$ \underline{by} $\mathbb G_{mS}$ \underline{in (\SGARefLink{sga7.sga7i.eix.n1-4-4}{1.4.4}),
so that $W_s$ is a biextension of $(A_o^o,{A'}_o^o)$ by}
$\mathbb G_{mk}$.  \underline{This biextension is separating}
(VIII~\SGARefLink{sga7.sga7i.eviii.n4-11}{4.11}).

Indeed, choose an ample invertible sheaf $\mathbb L$ on $A^o$.  Then
$\mathbb L_\eta$ defines a homomorphism (cf. \SGARefLink{sga7.sga7i.eix.n1-0-5}{1.0.5})
\[
 A_K\longrightarrow A'_K,
\]
which extends to a homomorphism $A\longrightarrow A'$, and hence gives
\[
 \psi:A^o\longrightarrow A'{}^o.
\]
We may consider the biextension $W_{\mathbb L}$ of $(A^o,A^o)$ by
$\mathbb G_{mS}$ obtained by pulling back $W$ along
$(\operatorname{id}_{A^o},\psi)$.  Its restriction to the generic
fibers is precisely the biextension $\delta(\mathbb L_\eta)$ of
VIII~\SGARefLink{sga7.sga7i.eviii.n4-12}{4.12}.  Hence, by VIII~\SGARefLink{sga7.sga7i.eviii.n7-4}{7.4}, $W_{\mathbb L}$ is associated with the
birigidified invertible sheaf $\delta(\mathbb L)$, and consequently
$(W_{\mathbb L})_s$ is defined by $\delta(\mathbb L_s)$.  On the other
hand, it may also be described as the pullback of $W_s$ by the
morphisms
\[
 \operatorname{id}_{A_o^o}:A_o^o\longrightarrow A_o^o,
 \qquad
 \psi_o:A_o^o\longrightarrow {A'}_o^o.
\]
Since $\mathbb L$, and hence $\mathbb L_s$, is ample, VIII~\SGARefLink{sga7.sga7i.eviii.n4-13}{4.13} shows
that $W_{\mathbb L_s}$ is separating, in particular separating on the
left.  A fortiori $W_s$ is therefore separating on the left.  By
symmetry it is also separating on the right, which proves \SGARefLink{sga7.sga7i.eix.n1-5}{1.5}.

We shall substantially sharpen \SGARefLink{sga7.sga7i.eix.n1-5}{1.5} when $A_o^o$ has unipotent rank
zero (\SGARefLink{sga7.sga7i.eix.n5-4}{5.4}).
